\chapter{Security Analysis Tables}\label{app:security}

This appendix gives the full security-analysis mapping tables summarized in
\cref{sec:eval-security}: the MITRE ATT\&CK technique-to-mitigation matrix
(\cref{sec:sec-mitre}) and the complete STRIDE enumeration (\cref{sec:sec-stride}). Both are drawn
from the attack scenarios of \cref{ch:threat} and the mechanisms of \cref{ch:design}, and both
correspond to the runnable adversarial suite of \cref{sec:impl-tests}. The identifiers are real
MITRE ATT\&CK Enterprise technique IDs; the tables map each to how it manifests against a key
directory and the key-transparency control that answers it.

% ----------------------------------------------------------------------------
\section{MITRE ATT\&CK Technique-to-Mitigation Matrix}\label{sec:sec-mitre}

The adversary techniques relevant to a messaging key directory are mapped to their
key-transparency mitigations in \cref{tab:sec-mitre}.

\begin{table}[ht]
\centering
\caption[MITRE ATT\&CK techniques mapped to KT mitigations.]{MITRE ATT\&CK Enterprise techniques mapped to their key-transparency mitigations. Two
technique families are mitigated: T1557 (with sub-technique T1557.004) and T1565 (with
sub-technique T1565.002); T1552.004 is the explicit endpoint-compromise boundary (X1).}
\label{tab:sec-mitre}
\footnotesize
\begin{tabular}{@{}p{1.7cm}p{2.2cm}p{3.1cm}p{4.5cm}l@{}}
\toprule
ID & Technique & Vector in KT & KT mitigation & Scen. \\
\midrule
T1557 & Adversary-in-the-Middle & network interception of X3DH & label mismatch (different secret $\Rightarrow$ different label) & A3 \\
T1557.004 & Evil Twin & rogue Wi-Fi access point for bundle interception & shared-secret divergence & A3 \\
T1565 & Data Manipulation & modify key bindings in the tree & Merkle inclusion-proof verification & A1,A3,A4 \\
T1565.002 & Transmitted Data Manipulation & replace a key in transit & self-check + warning-label propagation & A1 \\
T1552.004 & Private Keys & extract insecurely stored private keys from a device & \emph{none}: KT boundary (endpoint compromise) & X1 \\
\bottomrule
\end{tabular}
\end{table}

The techniques span the Collection, Credential Access, and Impact tactics. Every
technique except T1552.004 is mitigated by a key-transparency control; T1552.004 is the deliberate
endpoint-compromise boundary, out of scope for any directory mechanism (\cref{ch:threat}). The table
lists only techniques with a faithful ATT\&CK counterpart~\cite{mitre-attack}. Adversary behaviours
specific to a malicious \emph{directory} have no precise Enterprise-technique match and
are the thesis's own A1/A2/A4 threats, analyzed through the STRIDE enumeration and the attack
scenarios of \cref{ch:threat} instead. These behaviours include observing registration patterns,
operating a parallel encrypted channel, abusing the server's own write authority, selectively
censoring writes, and presenting divergent signed views (equivocation).

% ----------------------------------------------------------------------------
\section{STRIDE Enumeration}\label{sec:sec-stride}

The STRIDE enumeration covers the server (Merkle log, prefix tree, signed tree head), the client
(registration, verification, warning monitor), and the auditor. Of twenty threats, sixteen are
mitigated and four are residual; the summary is \cref{tab:eval-stride} in \cref{ch:eval}, and
the per-category detail follows in \cref{tab:stride-s,tab:stride-t,tab:stride-r,tab:stride-i,tab:stride-d,tab:stride-e}.

A threat is counted \emph{mitigated} when an implemented control prevents or detects the principal
attack outcome, even where a limitation remains; it is counted \emph{residual} when no implemented
control addresses that outcome. The Residual column records both cases: the four counted residual
threats are \emph{italicized}, while upright entries are limitations of rows that are counted
mitigated.

\begin{table}[ht]
\centering
\caption{Spoofing (4/4 mitigated).}
\label{tab:stride-s}
\footnotesize
\begin{tabular}{@{}p{4.4cm}p{5.2cm}p{4.0cm}@{}}
\toprule
Threat & KT control & Residual \\
\midrule
Forge a VRF label & ECVRF-EDWARDS25519-SHA512-TAI binding & none (cryptographic) \\
Impersonate a user registration & detected by the owner's self-check + peer warning (A1) & write API unauthenticated in the PoC; production requires account-layer authorization (\cref{ch:discussion}) \\
Forge a Merkle proof & SHA-256 collision resistance & negligible ($2^{128}$) \\
Forge an STH signature & Ed25519 verification & none (cryptographic) \\
\bottomrule
\end{tabular}
\end{table}

\begin{table}[ht]
\centering
\caption{Tampering (3/3 mitigated).}
\label{tab:stride-t}
\footnotesize
\begin{tabular}{@{}p{4.4cm}p{5.2cm}p{4.0cm}@{}}
\toprule
Threat & KT control & Residual \\
\midrule
Modify a key at a label & owner's self-check detects the mismatch & detection latency (1--2 epochs) \\
Alter a historical entry & append-only log + consistency proofs & none (cryptographic) \\
Tamper with a proof path & client-side \texttt{verify\_search\_result} & none (cryptographic) \\
\bottomrule
\end{tabular}
\end{table}

\begin{table}[ht]
\centering
\caption{Repudiation (3/3 mitigated).}
\label{tab:stride-r}
\footnotesize
\begin{tabular}{@{}p{4.4cm}p{5.2cm}p{4.0cm}@{}}
\toprule
Threat & KT control & Residual \\
\midrule
Deny a key change & append-only log records every mutation & none (immutable log) \\
Deny an attack occurred & warning entry + inclusion proof & none (cryptographic evidence) \\
Deny STH issuance & Ed25519 signature is non-repudiable & none (cryptographic) \\
\bottomrule
\end{tabular}
\end{table}

\begin{table}[ht]
\centering
\caption{Information Disclosure (2/3, 1 residual).}
\label{tab:stride-i}
\footnotesize
\begin{tabular}{@{}p{4.4cm}p{5.2cm}p{4.0cm}@{}}
\toprule
Threat & KT control & Residual \\
\midrule
Link a label to a user identity & VRF-derived pseudorandom labels & \emph{server knows the mapping (semi-trusted)} \\
Enumerate registered users & 256-bit pseudorandom labels & none for outsiders (operator residual counted above) \\
Leak the shared secret from a label & HKDF-Expand (dual-index labels) and ECVRF (identity labels) are one-way & none (cryptographic) \\
\bottomrule
\end{tabular}
\end{table}

\begin{table}[ht]
\centering
\caption{Denial of Service (1/4, 3 residual).}
\label{tab:stride-d}
\footnotesize
\begin{tabular}{@{}p{4.4cm}p{5.2cm}p{4.0cm}@{}}
\toprule
Threat & KT control & Residual \\
\midrule
Withhold epoch publication & auditor observes the head not advancing & \emph{no staleness timeout implemented; freshness alarm left to deployment} \\
Refuse to serve proofs & cross-vantage comparison by the auditor & \emph{gossip transport not implemented in the PoC; prototype fetches a second vantage directly} \\
Censor warning writes & compound signal (A2): key mismatch from the self-check + proof-verified \emph{absence} of the warning $=$ confirmed compromise & client out-of-band alert \\
Flood with fake registrations & deployment-level rate limiting & \emph{not implemented in the PoC (operational risk)} \\
\bottomrule
\end{tabular}
\end{table}

\begin{table}[ht]
\centering
\caption{Elevation of Privilege (3/3 mitigated).}
\label{tab:stride-e}
\footnotesize
\begin{tabular}{@{}p{4.4cm}p{5.2cm}p{4.0cm}@{}}
\toprule
Threat & KT control & Residual \\
\midrule
Server forges an inclusion proof & SHA-256 cryptographic binding & none (cryptographic) \\
Server fabricates a tree head & consistency-proof chain + Ed25519 STH & none (cryptographic) \\
Client escalates to server writes & write-once policy rejects overwrites of occupied labels & authenticated writes left to deployment (empty-slot squatting; \cref{ch:discussion}) \\
\bottomrule
\end{tabular}
\end{table}

The four residual risks are of two kinds, and none is a cryptographic weakness. One is inherent to the
semi-trusted-server model: the operator necessarily learns the label-to-user mapping (Information
Disclosure), a residual shared by the standard key-transparency literature (CONIKS~\cite{melara2015coniks},
the IETF Key Transparency architecture~\cite{ietf-keytrans-protocol}) and not addressed by
ACKA~\cite{dowling2023acka}, which abstracts the log behind a security definition rather than
realizing a distrusted one. The other three are availability functions the proof-of-concept does not
implement: write rate-limiting, a published-epoch staleness timeout, and a gossip transport for
proof-omission. Each is a deployment-layer engineering task rather than a gap in the scheme
(\cref{ch:discussion}).